\begin{frame}{현실 격차의 시각화}
  \begin{columns}[T]
    \begin{column}{0.62\textwidth}
      \centering
      \includegraphics[width=\linewidth]{ch13_1_domain_randomization.png}
      \vskip0.4em
      {\footnotesize 왼쪽: 20개 PD 정책을 감쇠 $\mu \in \{0.1,1,2,5\}$ 환경에서
      실행한 200스텝 누적 보상 곡선 (색 = $k_p$).
      오른쪽: 시뮬레이션($\mu = 0.1$) vs 실물($\mu = 5.0$) 산점도.}
    \end{column}
    \begin{column}{0.38\textwidth}
      \begin{itemize}
        \item $\mu$ 가 커질수록 점수는 낮아지되, $k_p = 8$ 계열이 가장
          \textbf{늦게} 떨어짐.
        \item $\mu = 5.0$ 에서는 $k_d$ 의 차이가 거의 사라짐.
        \item 산점도는 $y = x$ 대각선 주변에 모이나, \textbf{정책마다 다른
          편이}를 보임 --- 이 편이가 바로 ``현실 격차''.
      \end{itemize}
    \end{column}
  \end{columns}
\end{frame}
